% !PRINTED 335
\weberpage{335}
%<<CONT
racy in advance than the decrease of the differences
$\varDelta, \varDelta', \varDelta''\ldots$ If $\varDelta'$ is so small
that it falls outside the limits of the intended accuracy, then taking
the first difference into account gives an exact result. The rule for
the computation of roots set forth here is called the \gesperrt{Regula falsi}.

The simplest form of this rule, suited to the first approximation, is
this:

If a root $x$ of the equation $f(x)=0$ lies between $\alpha$ and
$\beta$, and if $f(\alpha) = -a, f(\beta) = b$, then
\weq{9}{ x = \frac{b\alpha + a\beta}{a+b} }
is an approximate value of $x$. This is only another way of writing
formula (7).

\begin{center}
§. 109.\\
\gesperrt{The Newton method of approximation.}
\end{center}

A method which, for the approximate computation of the roots of an
equation, is usually better than interpolation is due to
\gesperrt{Newton} and was developed and more closely investigated by
\gesperrt{Fourier}.

The method consists simply in the following. Let
\weq{1}{ f(x) = 0 }
be the equation to be solved, and let a value $x = \alpha$ have been
found which may be regarded as a certain approximation to a root. We
set in (1)
\[ x = \alpha + h, \]
and obtain
\weq{2}{ f(\alpha+h) = f(\alpha) + hf'(\alpha) +
  \frac{h^2}{1.2}f''(\alpha) + \cdots }
If one now determines $h$ from the equation
\weq{3}{ f(\alpha) + hf'(\alpha) = 0, }
which presupposes that $f'(\alpha)$ is different from zero, then
\[ f(\alpha+h) = \frac{h^2}{1.2}f''(\alpha) + \cdots, \]
and thus becomes, if $h$ is a small number, since only the square
%CONT>>
